\chapter{Implementation}\label{chap:chap4}
% This chapter presents the detailed implementation of MIRVisualScore, a modular Python-based visualization tool for music information retrieval. 
% Focus on design decisions, technical architecture, and practical usage.

\section{System Architecture and Design}

% High-level overview of MIRVisualScore's modular architecture
% Explain design philosophy: Layer-based composition pattern and Key design principles (modularity, extensibility, reusability)
% Explain the core abstraction: Layer class
% - Motivation for layer-based design
% Benefits: easy to compose, add new visualizations, maintain code
% - Show the base Layer class and its interface (load_data(), draw())
% - Explain how layers communicate through shared_data
% - Benefits for extensibility
% Present a system diagram showing component relationships

\section{Technology Stack and Dependencies}

% Subsection: Python and Core Libraries
% Describe versions and justification

% - Python 3.12.3+ (why this version)
% - Key dependencies and their purposes:
%   * matplotlib: visualization rendering
%   * numpy: numerical operations
%   * librosa or similar: audio processing
%   * Verovio or MEI tools: score parsing
%   * ScoreWarp: score-to-audio alignment
% - Why these technologies were chosen vs alternatives
% - Dependency management and requirements.txt


%?????????????
% \section{Core Modules}
% Brief description of the main porpuse 
% Beat_Layers
% Spectrogram_Layer
% visualization_system
% warp_score
%?????????????

\subsection{Data Loading and File Management}

% What file types are supported
% How data is loaded and cached
% Error handling for missing/invalid files
% File paths and input/output directory structure

% Describe LoadFiles class
% - Supported input formats: MEI scores, audio files (WAV, MP3), beat analysis files (NPZ)
% - Output formats and directory structure
% - File validation and error handling


\subsection{Visualization System Core}
% Visualizer class is the main orchestrator.
%
% Describe the Visualizer class in detail
% - Main entry point for creating visualizations
% - Layer composition and rendering pipeline: add_layer(), draw(), show()
%   - Methods for rendering: draw(), show()
% - Output generation: PNG and SVG export methods
% Data sharing between layers via shared_data dictionary.
% Matplotlib figure configuration
% - Configuration options: figsize, DPI, colors 

\subsection{Spectrogram Layer}

% Describe spectrogram computation and visualization
% - Audio loading and processing
% - Integration with audio timeline for synchronization
% Customization options for spectrogram display
% - STFT parameters (window size, hop length, etc.)
% - Frequency scale options (linear, mel, log)
% - Colormap selection and intensity normalization
% - Time-frequency resolution trade-offs

% Subsection 4.3.4: Beat Detection and Visualization (Beat_Layers.py)
% - Integration with BeatThis algorithm
% - Beat probability visualization
% - Onset detection and display
% - Temporal alignment with spectrogram and score
% - Customization: beat marker styles, colors, transparency

\subsection{Beat Detection and Visualization Layers}

% Describe beat tracking integration
% - BeatThis algorithm integration
% - Beat probability curves
% - Onset detection visualization
% - Beat confidence and annotation features
% - Performance considerations for real-time vs batch processing

\subsection{Score Alignment and Warping}

% Describe score-audio synchronization
% - Score file formats (MEI) and parsing
% - Onset detection for score notes (via MAPS JSON)
% - Score warping algorithm: ScoreWarp integration
% - Handling score timing variations
% - Visual representation of aligned score

\section{Data Pipeline and Processing Flow}

% Describe the full workflow
% - Input files: audio, scores, beat analyses
% - Processing stages: load → analyze → compose → render → export
% - Shared data dictionary structure
% - Synchronization mechanism between layers: establishing common time references
% - Output generation pipeline

\section{How to Use MIRVisualScore}

% Subsection: Installation and Setup
% - Virtual environment creation
% - Dependency installation
% - Verification steps

% How to get the required files:
% 1) How to get the Aligned Score
% 2) How to get/generate the MAPS file
%       Autogenerate Maps file  and Maps file structure
% 3) How to get .npz

% Creating a basic visualization (explain example.py)
%   - Loading Audio
%   - Adding multiple layers (onsets, beats, downbeats etc)
%   - Adding the Score

\subsection{Installation and Setup}

% Step-by-step setup guide

% Subsection: Basic Usage Example
% - Simple "hello world" visualization example
% - Loading audio file
% - Adding spectrogram layer
% - Exporting result

\subsection{Basic Usage Example}

% TODO: Provide minimal working example with code

% Subsection: Advanced Examples
% - Combining multiple layers (spectrogram + beats + score)
% - Using different customization options
% - Working with different audio formats and scores

\subsection{Advanced Usage Examples}

% Provide examples of real-world use cases


\section{How to Further Customize Visualizations}

% Subsection: Creating Custom Layers


\subsection{Creating Custom Visualization Layers}

% Guide for extending MIRVisualScore with custom layers

% Subsection: Styling and Customization
% - Color palettes and colormaps
% - Figure layout customization
% - Font sizes and labels
% - Axis ranges and scales

\subsection{Styling and Aesthetic Customization}

% Describe customization options for appearance

% Subsection: Integration with External Libraries
% - Combining with other MIR tools
% - Output processing pipelines
% - Web deployment considerations

\subsection{Integration and Interoperability}

%  Discuss how to integrate with other systems

\section{Implementation Challenges and Design Decisions}

% Discuss challenges in:
% - Key technical challenges encountered during development
%      Audio-score synchronization.
%       ???????

\section{Summary}

% TODO: Summarize key aspects of the implementation
% - Modular architecture enables extensibility
% - Layer-based design provides flexibility
% - Integration of multiple visualization types
% - Balance between ease-of-use and advanced customization
% - Contribution to MIR research tools


